\section{Supplementary Methods}
\label{sec:SI_methods}
Here I provide the technical details of the procedure summarized in Methods.

\renewcommand{\thesubsection}{S\arabic{subsection}}
\setcounter{subsection}{0}

\subsection{Derivation of PCA and tICA}


\label{subsec:SI_dimred}
The presentation of tICA in this section follows closely that of \cite{my_tica_ref}.


PCA and tICA are both linear dimensionality reduction techniques, i.e. they project the data onto a linear subspace of the original feature space. 
Let $\{\mathbf{x}(t)\}_{t=0}^{T-1}$ be a multidimensional, discrete time-series, where each snapshot is a column vector, 
of dimension $D$, corresponding to an arbitrary, vectorized representation of a protein conformation (e.g. a set of dihedral angles). 
Suppose the data has been sampled from a \textit{stationary} and \textit{time-reversible} stochastic process. 
Without loss of generality, I assume that the data have been centered, $\mathbb E[\mathbf{x}] = 0 $. Using the braket notation $\bra{x} := \mathbf{x}^T, \ket{x} := \mathbf{x} $, I define the covariance matrix:


\begin{equation}
\mathbf{C}_0 := \mathbb{E}\!\left[\ket{x}\bra{x}\right]
\end{equation}

and the time-lagged covariance matrix:
\begin{align}
    \mathbf{C}_{0,\tau}:&= \mathbb{E}\left[\ket{x(t)}\bra{x(t + \tau)}\right].
\end{align}

For a stationary reversible process,
\begin{equation*}
\mathbf C_{0,\tau} = \mathbf C_{\tau,0}^{T} = \mathbf C_{\tau,0},
\end{equation*}
so that the lagged covariance matrix may be denoted simply by \(\mathbf C_\tau\).

If $\ket{\hat{u}}\bra{\hat{u}}$ is a projection operator (i.e. $\ket{\hat{u}}$ is a unit vector), then:

\begin{equation}
    \bra{\hat{u}}\mathbf{C}_0\ket{\hat{u}} = \mathbb{E}[\braket{\hat{u}|x}\braket{x|\hat{u}}] = \mathbb{E}[\braket{\hat{u}|x}^2]
\end{equation}
is the variance of the projection of the data along $\ket{\hat{u}}$, and

\begin{align}
    \label{eq:autocorrelation}
    \textrm{ACF}_{\hat{u}}(\tau) &=\frac{\bra{\hat{u}}\mathbf{C}_{\tau}\ket{\hat{u}}}{\bra{\hat{u}}\mathbf{C}_{0}\ket{\hat{u}}} \notag \\
    \\ &= \frac{\mathbb{E}[\braket{\hat{u}|x(t)}\braket{\hat{u}|x(t+\tau)}]}{\mathbb{E}[\braket{\hat{u}|x}^2]} \notag 
\end{align}

its autocorrelation at lag time $\tau$.

\vspace{0.5em}

\noindent \textbf{PCA} seeks the direction of maximal variance. The first principal component is defined as
\begin{align}
\ket{\hat{u}_1} =& \arg\max \bra{\hat{u}}\mathbf{C}_0\ket{\hat{u}}  \notag \\
&\text{subject to}\,\,\|\hat{u}\|=1
\end{align}

this optimization problem can be solved introducing a Lagrange multiplier $\lambda$, and the stationarity condition yields

\begin{equation}
\mathbf{C}_0 \ket{\hat{u}_i} =\lambda_i \ket{\hat{u}_i},
\end{equation}

i.e. the projection versor is an eigenvector of the covariance matrix. The corresponding eigenvalue equals the variance of the data projected onto the principal component:

\begin{equation*}
\bra{\hat{u}_i}\mathbf{C}_0\ket{\hat{u}_i} = \lambda_i \cdot \braket{\hat{u}_i|\hat{u}_i} = \lambda_i.
\end{equation*}

Then the second principal component is defined as the direction of maximal variance orthogonal to the first, giving the 
optimization problem:
\begin{align}
\ket{\hat{u}_2} = &\text{argmax}\bra{u}\mathbf{C}_0\ket{u} \notag \\
&\text{subject to}  \,\, \braket{u|u} = 1 \notag \\
&\text{and}   \,\, \braket{u|\hat{u}_1} = 0
\end{align}

which again can be solved with Lagrange multipliers. Iterating, this implies that the principal components form an orthonormal basis of eigenvectors of the covariance matrix, i.e. they diagonalize the covariance matrix.
Dimensionality reduction ($d < D$) is achieved by taking the components along the first $d$ 
principal components $\{\ket{\hat{u}_1}, \ldots, \ket{\hat{u}_d}\}$.


\begin{equation}
z_i = \braket{\hat{u}_i|x}, \qquad \mathbf z = U^T \mathbf x.
\end{equation}

This is a linear projection onto an ortonormal basis:

\begin{equation*}
    \ket{x} \rightarrow \ket{z} := \sum_{i = 1}^{d}\, \braket{x | \hat{u}_i}\, \ket{\hat{u}_i} 
\end{equation*}

By construction, in the reduced space components are uncorrelated and ordered by decreasing variance:

\begin{equation}
 \mathbb{E}\!\left[\ket{z}\bra{z}\right] = \text{diag}(\lambda_1,\, \cdots \lambda_d)
\end{equation}
\vspace{1em}

\noindent \textbf{tICA} Unlike PCA, tICA \cite{my_tica_ref} seeks directions $\ket{u}$ that maximize the autocorrelation Eq.\eqref{eq:autocorrelation} at a fixed lag time $\tau$. The lagtime $\tau$ is a free parameter.
The autocorrelation is invariant under scaling of the projection vector $\ket{u}$. Differently from PCA, the normalization constraint is usually implemented as $\bra{u}\mathbf{C}_{0}\ket{u}=1$, as this choice simplifies the optimization problem.
Formally, tICA solves the optimization problem:

\begin{align}
    \label{eq:generalized_eigenvalue_problem}
\ket{\hat{u}_1} =& \text{argmax}\bra{u}\mathbf{C}_{\tau}\ket{u} \notag \\
&\text{subject to}  \,\, \bra{u}\mathbf{C}_{0}\ket{u} = 1
\end{align}

using the same procedure as for PCA, I can solve this with Lagrange multipliers, and show that


\begin{equation}
\mathbf{C}_{\tau}\ket{\hat{u}_1}
=
\lambda_1
\mathbf{C}_{0}\ket{\hat{u}_1},
\end{equation}.

The generalized eigenvalue $\lambda_i$ is equal to the autocorrelation of the projected coordinate $\braket{\hat{u}_i|x(t)}$ at lag time $\tau$. Assuming that this coordinate approximates a single relaxation mode of the dynamics, its autocorrelation decays exponentially, yielding

\begin{equation}
    \label{eq:eigvals_and_timescales}
    \lambda_i \approx e^{-\tau/t_i},
\end{equation}

where $t_i$ is the corresponding relaxation timescale. In practice, the relation is approximate because a tIC may contain contributions from multiple dynamical modes.

Similarly as done with PCA, I can iterate this procedure and define the second tIC as the direction of maximal autocorrelation at lag time $\tau$ that is \textit{uncorrelated} with the first:


\begin{align}
\ket{\hat{u}_2} = &\text{argmax} \bra{u}\mathbf{C}_{\tau}\ket{u} \notag \\
&\text{subject to}  \,\, \bra{u}\mathbf{C}_{0}\ket{u} = 1 \notag \\
&\text{and}   \,\, \bra{u}\mathbf{C}_{0}\ket{\hat{u}_1} = 0
\end{align}

the solution $\ket{\hat{u}_1}$ is again a solution of the same generalized eigenvalue problem as in Eq.\eqref{eq:generalized_eigenvalue_problem}.
generalized eigenvalue problem.
Dimensionality reduction ($d < D$) is achieved by taking the components of the data $\mathbf{x}(t)$ along the first $d$ tICs $\{\ket{\hat{u}_1}, \ldots, \ket{\hat{u}_d}\}$.

\begin{equation}
    \mathbf{z}(t) = (\braket{x(t) | \hat{u}_i}, \cdots, \braket{x(t) | \hat{u}_d})
\end{equation}

By construction, in the reduced space components $z_i$ have unit variance, are uncorrelated at zero lag,  and are ordered by decreasing autocorrelation at lag time $\tau$.

\begin{equation}
    \label{eq:original_tica}
 \mathbb{E}\!\left[\ket{z}\bra{z}\right] = \mathbb{I}
\end{equation}

\vspace{1em}


\noindent \textbf{tICA as a diagonalization in the whitened coordinate space}
It is worth noting that unlike in PCA, the tIC vectors are generally not orthogonal under the standard Euclidean inner product. Instead, as I said, they satisfy the generalized orthogonality relation
$$
\bra{u_i}\mathbf C_0\ket{u_j}=\delta_{ij}.
$$
i.e. they corresponding components $z_i$ and $z_j$ are \textit{uncorrelated} at zero lag. Nevertheless, the tICA procedure can be reformulated as a standard eigenvalue problem in a transformed coordinate system. The key idea is to first apply a whitening transformation, which rescales and rotates the original data so that all directions have unit variance and are mutually uncorrelated. Define the ZCA whitening transformation

\begin{equation*}
\ket{y} = \mathbf C_0^{-1/2} \ket{x}
\end{equation*}

where $C_0^{-1/2}$ is symmetric and positive definite. By construction, the covariance matrix in the transformed space is the identity:

\begin{equation*}
\mathbb E\left[\ket{y}\bra{y}\right] = \mathbf{C}_0^{-1/2} \mathbf{C}_0 \mathbf{C}_0^{-1/2} = \mathbf I.
\end{equation*}

and time-lagged covariance matrix in the whitened space is

\begin{equation*}
\hat{\mathbf{C}}_\tau = \mathbf C_0^{-1/2} \mathbf C_\tau \mathbf C_0^{-1/2}
\end{equation*}

$\hat{\mathbf{C}}_\tau$ is symmetric and can thus be diagonalized:,

\begin{equation*}
\hat{\mathbf{C}}_\tau \ket{v_i} = \lambda_i \ket{v_i}, \qquad
\braket{v_i|v_j} = \delta_{ij}.
\end{equation*}

If $\ket{u_i}$ solves the tICA problem \eqref{eq:generalized_eigenvalue_problem}, then $\ket{v_i} = \mathbf C_0^{1/2} \ket{u_i}$ is an eigenvector of $\hat{\mathbf{C}}_\tau$ with the same eigenvalue.

Indeed, multiplying the generalized eigenvalue problem \eqref{eq:generalized_eigenvalue_problem} from the left by $\mathbf C_0^{-1/2}$ and defining
$\ket{v_i} = \mathbf C_0^{1/2}\ket{u_i}$ yields

\begin{equation*}
\hat{\mathbf C}_\tau \ket{v_i} = \lambda_i \ket{v_i}.
\end{equation*}

Consequently, the tICA components $z_i$ correspond to the ortogonal projection of the whitened data $\mathbf{y}$ onto the ortonormal basis defined by $\{\ket{v_i}\}$:

\begin{equation*}
z_i(t) = \braket{x(t)|u_i} = \braket{x(t) | C_0^{-1/2} v_i}  = \braket{y(t) | v_i}.
\end{equation*}


\vspace{1em}
\noindent \textbf{Alternative generalized eigenvector normalizations for tICA}\newline 
The construction presented above corresponds to the original formulation of tICA, in which the covariance matrix of the transformed data is the identity \eqref{eq:original_tica}.
HoIver, alternative rescaling of the transformed coordinates have been used in the literature and are usually implemented in softwares to endow Euclidean distances in the reduced space with a more direct kinetic interpretation. These can 
be particularly useful when the tICA coordinates are subsequently clustered to construct a Markov state model, as I do in this work.

The most commonly used choice is the \textit{kinetic map} in which each tIC component $z_i$ is multiplied by its corresponding eigenvalue,

\begin{equation}
\tilde z_i(t) = \lambda_i z_i(t).
\end{equation}

With this normalization, the coordinates representing the sloIst processes "Iight more", and a unit step in those directions in the reduced space is assigned a larger distance.
An alternative is the \textit{commute map}, where each coordinate is rescaled according to its implied timescale,
\begin{equation}
\tilde z_i(t) = \sqrt{\frac{t_i}{2}} z_i(t), \qquad t_i = -\frac{\tau}{\ln \lambda_i}.
\end{equation}
With this normalization, Euclidean distances approximate commute distances, i.e. the expected time required to travel from one configuration to another and back. This is a variation on the same idea behind the kinetic map, being eigenvalues and timescales connected by \eqref{eq:eigvals_and_timescales}.



\subsection{Elements of Markov theory}
\label{suppsec:markov_theory}

\textcolor{gray}{
Markovianity assessment and optimal lagtime choice was performed through the implied-timescale convergence test. The i-th implied timescale estimated at lag $\tau$,  $t_i(\tau)$ is given by the relation:
\begin{equation}
t_i(\tau) = -\frac{\tau}{\ln \lambda_i(\tau)}
\end{equation}
where $\lambda_i(\tau)$ is the $i$-th eigenvalue ($i \geq 2$) of the transition matrix estimated at lag time $\tau$.
It is known that the true timescales of the underlying continuous-state Markov process, $t_i^{*}$, are approximated from below ($t_i(\tau) \leq t_i^{*}$) and the relative error decreases as
i) the discretization error decreases or ii) the lagtime $\tau$ increases. With this consideration, the implied timescale test consists in plotting $t_i(\tau)$ as a function of $\tau$ and looking for a region of approximate convergence, where the estimated timescales are approximately constant.
The lagtime is selected as the minimum value in the region of convergence, as this choice minimizes the loss of temporal resolution while ensuring that the process of interest is approximately Markovian \cite[Ch. 4.7]{thebible}. 
}

\textcolor{gray}{
After selecting the MSM lag time based on the implied timescale convergence analysis, model quality was further assessed using an autocorrelation function test. 
While a direct Chapman-Kolmogorov (CK) test provides a stringent validation of Markovianity, its application at the microstate level becomes impractical for models containing hundreds of states, 
as it requires comparing $T(k\tau)$ and $[T(\tau)]^k$ entry-wise, and summary statistics are needed which could be difficult to interpret. 
Instead, following the approach suggested by [Chap. 2.4.2]\cite{thebible}, model predictions were evaluated through the autocorrelation function of a physically meaningful observable.
As an observable, the root-mean-square deviation (RMSD) from the crystal structure was chosen. For each microstate, a representative observable value was obtained by averaging the RMSD over all trajectory
frames assigned to that microstate. The MSM prediction for the autocorrelation function was then computed from the spectral decomposition of the transition matrix,
}

\noindent The dynamics encoded by a Markov state model can be understood through the eigendecomposition of its transition matrix $\mathbf{T}$.
For consistency with the notation used in the Python package deeptime, let us consider the transition matrix $\mathbf{T}$ as a row-stochastic matrix, i.e. $T_{ij}$ is the probability of transitioning from state $i$ to state $j$ in one lag step:
\begin{equation*}
    \sum_j T_{ij} = 1, \qquad T_{ij} \geq 0.
\end{equation*}
Probability distributions are represented as row vectors, and evolve as:
\begin{equation*}
\mathbf{p}(n) = \mathbf{p}(0)\mathbf{T}^n,
\end{equation*}
where $\mathbf{p}(0)$ is the initial distribution and $\mathbf{p}(n)$ is the distribution after $n$ lag steps.
The transition matrix admits a spectral decomposition of the kind:
\begin{equation*}
\mathbf{T} = \sum_i \lambda_i \mathbf{r}_i \mathbf{l}_i^{\mathrm T},
\end{equation*}
where $\lambda_i$ are the eigenvalues, and $\mathbf{r}_i$ and $\mathbf{l}_i$ are the corresponding right and left eigenvectors, satisfying:
\begin{equation*}
\mathbf{T}\mathbf{r}_i=\lambda_i\mathbf{r}_i,
\qquad
\mathbf{l}_i^{\mathrm T}\mathbf{T}=\lambda_i\mathbf{l}_i^{\mathrm T}.
\end{equation*}


The probability distribution after $n$ lag times can therefore be written as

\begin{equation*}
\mathbf{p}(n)=
\mathbf{p}(0)\mathbf{T}^{,n}=
\sum_i
c_i\, \cdot \lambda_i^n \cdot \mathbf{l}_i,
\end{equation*}

where the coefficients $c_i = \mathbf{p}(0)\mathbf{r}_i$ depend on the initial condition. The dominant eigenvalue is $\lambda_1=1$ and corresponds to the equilibrium distribution. All other modes decay exponentially with characteristic relaxation times

\begin{equation*}
t_i
-\frac{\tau}{\ln \lambda_i},
\end{equation*}

where $\tau$ is the MSM lag time.

The left eigenvectors therefore describe the spatial structure of the relaxation processes occurring on the corresponding timescales. States with eigenvector components of opposite sign exchange probability during the relaxation process, while the magnitude of each component quantifies the degree to which that state participates in the process. The instantaneous probability flux associated with a mode can be visualized from the change in the probability distribution after one lag step,

\begin{equation*}
\Delta\mathbf{p} = 
\mathbf{p}\mathbf{T} - \mathbf{p} = {\color{orange} TODO: complete}
\end{equation*}

For reversible Markov models obeying detailed balance, left and right eigenvectors contain equivalent information and are related through the equilibrium population matrix $\boldsymbol{\Pi}$,

\begin{equation*}
\mathbf{l}_i
\boldsymbol{\Pi},\mathbf{r}_i .
\end{equation*}

Consequently, the right eigenvectors may be interpreted as equilibrium-population-weighted versions of the relaxation modes.



Autocorrelation of an observable $\theta$ can be expressed in terms of the spectral decomposition of the transition matrix. Let $\theta_i$ be the value of the observable in microstate $i$, and let $\langle \theta, \mathbf{r}_i\rangle_\pi$ denote the inner product weighted by the stationary distribution $\pi$. The autocorrelation function at lag time $n$ is then given by

\begin{align}
c(n)
&=
\frac{
\sum_{i=1}^{N}
\lambda_i^{n}\,
\langle \theta,\mathbf{r}_i\rangle_\pi\,
\langle \theta,\mathbf{l}_i\rangle
}
{
\sum_{i=1}^{N}
\langle \theta,\mathbf{r}_i\rangle_\pi\,
\langle \theta,\mathbf{l}_i\rangle
}
\notag \\
&\equiv
\frac{
\sum_{i=1}^{N}
\lambda_i^{n}\,
\langle \theta,\mathbf{l}_i\rangle^2
}
{\sum_{i=1}^{N}\, \langle \theta,\mathbf{l}_i\rangle^2} \notag
\end{align}

where $\lambda_i$ and $(\mathbf{l}_i, \mathbf{r}_i)$ denote the eigenvalues and (left, right) eigenvectors of the transition matrix (in row-stochastic notation), $\pi$ the stationary distribution,
and ($\boldsymbol{\theta}$) is the vector of microstate-averaged RMSD values. Equivalence between the two reported formula hold because of the detailed balance assumption.
%


\textcolor{gray}{
The dynamical processes captured by each MSM were characterized through the spectral decomposition of the transition matrix (see Supplementary Information \ref{si:spectral} for details). 
The eigenvalues were converted into relaxation timescales, while the corresponding eigenvectors were used to identify the states participating in each relaxation process. 
States with eigenvector components of opposite sign exchange probability on the associated timescale, whereas the magnitude of each component reflects the degree of participation
of that state in the process. For reversible models satisfying detailed balance, left and right eigenvectors contain equivalent kinetic information
and are related through the equilibrium distribution.
The signs of eigenvectors components provide a criterion for defining partitions in the microstate model. This is systematically done through the algorithm PCCA++ \cite[Ch. 2.5.1 and 2.5.2]{thebible}. I use this method to connect
the slowest exchange modes in the markov chain to the thermodynamics of the protein, although the two are not strictly the same thing. 
}
\vspace{0.1em} \newline \textit{Lagtime choice:}
Estimation of the transition matrix requires specification of the lag time $\tau$ ($= 1, 2, 3\cdots$ in number of frames) at which to count transitions between states. Choice of lag time is done through the
timescale convergence test. Briefly (see Supplementary Methods \ref{suppsec:markov_theory} for more details), the transition matrix $T(\tau)$  is estimated at different lag times and its eigenvalues are used to compute the implied timescales $t_i(\tau)$, which are then plotted as a function of $\tau$. 
If the Markovian assumption is valid, the implied timescales should be approximately constant for sufficiently large $\tau$. The lag time is then selected as the minimum value in the region of convergence, as this choice minimizes the loss of temporal resolution while ensuring that the process of interest is approximately Markovian\cite[Ch. 4.7]{thebible}.
\vspace{0.1em} \newline \textit{MSM validation:} to validate the model against the original trajectory, I compare the autocorrelation function of the root mean square deviation (RMSD) to the crystal structure computed from the trajectory (\texttt{hp35.crystaldist}) and predicted by the MSMs (see \ref{suppsec:markov_theory}).
%